\subsection{Providerneutrale Persistenzstrategie für Produkt- und Nutzerdaten}
\label{subsec:providerneutrale-persistenzstrategie}

Die Baseline trennt Template- und Entwicklungsdaten, Laufzeitkonfiguration
sowie Produkt- und Nutzerdaten. \texttt{profiles/},
\texttt{project-profile.toml}, \texttt{tools/} und \texttt{shared/}
beschreiben, erzeugen oder prüfen das Projekt. \texttt{.env},
\texttt{config/environment.toml}, \texttt{DATABASE\_URL} und
API-Konfiguration bestimmen dagegen, wie eine Instanz ausgeführt wird.
Dokumente, Planungsstände, fachliche Datensätze, Medien, Projektdateien und
Historien entstehen erst bei der Produktnutzung und benötigen einen eigenen
Verantwortungsbereich und Lebenszyklus. Konfigurationsdateien sind kein
Nutzerdatenspeicher; versionierte Template-Assets sind keine Laufzeit-Assets
des fertigen Produkts.

Eine Persistenzentscheidung unterscheidet providerneutral lokale
Dateispeicherung, eingebettete Datenbanken, entfernte Datenbanken und
Asset-Speicher. Nicht jedes Produkt benötigt alle vier Klassen. Lokale Dateien
können etwa JSON, Markdown oder produktspezifische Formate verwenden; eine
eingebettete Datenbank könnte beispielsweise SQLite nutzen. Die bestehenden
Capabilities \texttt{database} und \texttt{postgres} implementieren dagegen
serverseitige SQL-Persistenz mit SQLAlchemy, Alembic und PostgreSQL.
Asset-Speicher für Bilder, Dokumente, Grafiken oder Binärdateien bleibt von
strukturierten Daten getrennt. \texttt{local-files},
\texttt{embedded-database}, \texttt{sqlite}, \texttt{object-storage} und
Synchronisation sind mögliche Extension Paths, keine implementierten
Capabilities \parencite{kleiveist2026persistence}.

Jedes Produktprojekt muss die maßgebliche Datenquelle festlegen. Bei einer
lokalen Anwendung kann ein lokaler Speicher die Source of Truth sein, bei
einer serverzentrierten Anwendung eine zentrale Datenbank. Offline- oder
Local-first-Systeme benötigen zusätzlich explizite Regeln für
Synchronisationsrichtung, Revisionen und Konflikte; eine
Synchronisationsengine stellt die Baseline nicht bereit. Zur
Produktentscheidung gehören außerdem Schema- und Versionsstrategie,
Migration, Backup und Recovery sowie Sicherheitsgrenzen.

Das Template standardisiert somit nicht die Daten selbst, ein universelles
Format oder einen allgemeinen Storage Provider, sondern die
Verantwortungsgrenzen und die zu dokumentierenden Entscheidungen für ihren
Umgang. Speicherformat und Persistenzarchitektur, Plattformprofil und Storage
Provider, Konfiguration und Nutzerdaten sowie Asset-Speicher und strukturierter
Datenspeicher bleiben jeweils getrennte Entscheidungen.
\beginnerinput{workspace/statement/300_architektur/380}
